\chapter{绪论}

\section{研究背景与问题定义}

PostgreSQL 是当前最具代表性的开源关系型数据库管理系统之一，其代码规模大、模块耦合复杂、演化周期长，长期服务于金融、政企与互联网等高可靠场景。随着版本持续迭代，系统在性能、并发能力与功能丰富度不断提升的同时，也不可避免地产生与修复各类缺陷。相比单次 Bug 个案分析，面向长期提交历史的系统性研究更有助于回答两个核心问题：其一，Bug 修复在时间、作者与模块维度上是否存在稳定规律；其二，修复行为在代码结构层与逻辑约束层是否可被统一建模与验证。

现有工程实践中，对 Bug 修复质量的判断通常依赖代码审查与回归测试。然而，测试本质上只能证明“存在问题”或“在有限样例下未发现问题”，难以在全输入空间上证明修复逻辑的完备性。特别是“静默逻辑失效（Silent Failure）”这类非崩溃型缺陷，可能在返回码正常的情况下造成行为偏差，传统测试覆盖不足时容易长期潜伏。因此，本研究将静态结构分析、动态执行追踪与形式化验证联合起来，尝试建立一条可复现、可解释、可证明的 Bug 修复分析路径。

基于上述背景，本文将研究对象限定为 PostgreSQL 主分支近年提交，围绕“如何识别高置信 Bug 修复提交、如何抽象修复结构模式、如何验证关键修复的正确性”展开。研究问题可形式化为：
\begin{enumerate}
  \item Bug 修复提交在时间分布、作者贡献和模块热点上是否具有可量化规律；
  \item 修复前后代码在条件判断、控制流和错误处理路径上是否存在可迁移的结构演化模式；
  \item 能否将动态分析发现的逻辑漏洞抽象为约束模型，并利用求解器给出“修复前可构造反例、修复后不可构造反例”的严格证据。
\end{enumerate}

\section{研究目标与贡献}

围绕以上问题，本文设定如下研究目标：
\begin{itemize}
  \item 构建 PostgreSQL Bug 修复提交的识别与标注流程，形成可复查的数据基础；
  \item 通过静态分析提炼修复行为的结构化特征，给出可解释的模式归纳；
  \item 通过动态追踪复现关键案例，定位非崩溃型逻辑缺陷并验证修复效果；
  \item 通过 Z3 形式化建模，将“测试结论”提升为“逻辑正确性证明”。
\end{itemize}

本文的主要贡献体现在以下三个方面：
\begin{enumerate}
  \item \textbf{多证据链研究框架}：提出“提交统计--静态结构--动态轨迹--形式化验证”的闭环分析框架，兼顾宏观规律与微观机制；
  \item \textbf{结构化修复模式抽象}：归纳出前置条件显式化、错误路径显式化、条件语义增强等可迁移模式，为后续自动审计与自动修复提供结构基础；
  \item \textbf{可证明的修复有效性}：在典型配置同步案例上，以 SAT/UNSAT 对照方式证明修复前后逻辑差异，增强结论的严格性与工程可信度。
\end{enumerate}

\section{研究方法与技术路线}

本文采用分层方法设计，技术路线如图~\ref{fig:tech-route} 所示（可在后续章节补充正式图示）：
\begin{enumerate}
  \item \textbf{提交数据采集与预处理}：围绕目标时间段提取提交元数据与差异信息；
  \item \textbf{Bug 修复识别与标注}：结合关键词信号、结构信号与测试信号，对候选提交进行 Y/M/N 分级标注；
  \item \textbf{静态结构演化分析}：使用 AST/libcst 对修复前后代码结构进行对比，统计条件节点、分支路径与错误处理变化；
  \item \textbf{动态执行轨迹分析}：基于 PySnooper 与测试脚本复现实例，捕获变量演化与异常触发路径；
  \item \textbf{形式化建模与验证}：将关键逻辑抽象为状态与约束，利用 Z3 检验是否存在违反规约的反例。
\end{enumerate}

\begin{figure}[H]
  \centering
  % \includegraphics[width=0.85\textwidth]{figures/tech_route.pdf}
  \caption{本文技术路线示意图（占位符，后续替换为正式图）}
  \label{fig:tech-route}
\end{figure}

为确保结果可复现，本文在每个阶段都保留了对应产物：识别规则文档、标注流程说明、动态日志、测试结果及求解脚本。通过多源证据互证，尽量减少单一方法的偏差风险。

\section{论文结构}

全文共七章，章节安排如下：
\begin{itemize}
  \item \textbf{第一章 绪论}：说明研究背景、问题定义、目标贡献与技术路线；
  \item \textbf{第二章 数据集构建与 Bug 修复识别}：介绍数据来源、识别规则、标注流程与统计结果；
  \item \textbf{第三章 Bug 修复结构演化分析（静态）}：给出 AST/libcst 视角下的结构变化与模式归纳；
  \item \textbf{第四章 动态执行轨迹与案例复现}：展示关键案例的异常定位、修复策略与回归验证；
  \item \textbf{第五章 形式化建模与修复正确性验证}：给出 Z3 建模过程与 SAT/UNSAT 结果对照；
  \item \textbf{第六章 综合讨论}：分析三类证据的互补关系、工程启示与有效性威胁；
  \item \textbf{第七章 结论与未来工作}：总结研究发现并给出后续拓展方向。
\end{itemize}

\section{本章小结}

本章围绕 PostgreSQL Bug 修复研究的现实需求，明确了本文的问题边界、研究目标与贡献定位，并给出了覆盖“识别--分析--验证”的整体技术路线。后续章节将在统一数据与规则基础上，分别展开静态、动态与形式化三条分析主线，最终形成可复现且可证明的结论体系。
